\documentclass[../main.tex]{subfiles}
\graphicspath{{\subfix{../image/}}}

\begin{document}

\chapter{微分形式}\label{chapter:ch14-differential-forms}

本章是关于微分形式理论的引论，它在光滑流形上的微分几何、拓扑与微分方程中起着核心作用。它使我们能够把向量微积分中的许多构造推广到流形上，例如曲线与曲面上的向量场导数。

在本章中，我们主要关注微分形式的代数与微积分类性质。我们将要建立的理论将在第~15 章中扮演关键角色，在那里我们把它与流形上的定向概念联系起来；而在第~16 章中，我们将利用它来定义流形上的积分并证明 Stokes 定理。

我们从斜对称张量的代数开始，即从一个向量空间的对偶空间上某些满足斜对称性的多重线性泛函开始。我们将看到，这些张量具有丰富的代数结构，包括一个乘积运算，称为外积。然后我们把斜对称张量放到流形上，从而得到微分形式，并说明微分形式如何在局部坐标下表示。在最后一节中，我们引入外导数运算，它作用于微分形式，并把关于光滑函数微分的概念推广到任意次数的微分形式。

\subfile{Chapter-14/section-01.tex}
\subfile{Chapter-14/section-02.tex}
\subfile{Chapter-14/section-03.tex}
\subfile{Chapter-14/section-04.tex}

\end{document}
